\documentclass[11pt,a4paper]{article}

\usepackage{kotex}
\usepackage{fontspec}
\setmainfont{Times New Roman}
\setmainhangulfont{AppleMyungjo}
\setsanshangulfont{Apple SD Gothic Neo}

\usepackage[a4paper,top=18mm,bottom=18mm,left=20mm,right=20mm]{geometry}
\usepackage{array}
\usepackage{enumitem}
\linespread{1.18}
\setlength{\parindent}{0pt}

\newcommand{\sol}[2]{\vspace{6pt}{\sffamily\bfseries #1} \quad 정답 \fbox{#2}\par\vspace{1pt}}

\begin{document}

\begin{center}
{\fontsize{15}{17}\selectfont\sffamily\bfseries 2026 고1 영어 변형 — 지문 23 정답 및 해설}\\[2pt]
{\sffamily 죄책감과 수치심}
\end{center}
\vspace{6pt}

\begin{center}
\renewcommand{\arraystretch}{1.4}
\begin{tabular}{|c|c|c|c|c|}
\hline
\sffamily 문항 & 1 & 2 & 3 & 4 \\\hline
\sffamily 유형 & \sffamily 빈칸추론 & \sffamily 영어 제목 & \sffamily 요약문 (A)(B) & \sffamily 서술형 해석 \\\hline
\sffamily 정답 & ③ & ③ & ③ & (아래) \\\hline
\sffamily 배점 & 2점 & 2점 & 2점 & 3점 \\\hline
\end{tabular}
\end{center}
\vspace{8pt}\hrule\vspace{4pt}

\sol{1. 빈칸추론}{③}
빈칸은 ``\textit{shame often is \underline{\hspace{2cm}}}''에 해당한다. 지문은 guilt의 효과는 종종 긍정적(\textit{positive})인 반면 shame은 그 반대임을 대조하며, 마지막 문장에서 ``\textit{shame \textbf{corrodes} the part of us that believes we can change}''라고 명시한다. \textit{corrosive}(부식성의, 좀먹는)는 동사 \textit{corrode}와 직접 연결되며 shame의 파괴적 특성을 가장 정확히 표현한다. 따라서 정답은 ③ \textbf{corrosive}.
\par\vspace{3pt}
\textbf{오답 소거:} ① motivating — 지문에서 \textit{motivator}는 shame이 아닌 \textbf{guilt}의 특성으로 제시됨. ② constructive — 본문의 뜻(부정적)과 정반대. ④ unconscious — 지문에 근거 없음. ⑤ inevitable — 지문에 근거 없음.

\sol{2. 영어 제목}{③}
글의 핵심 주장: guilt는 긍정적 행동 변화를 이끄는 동기(\textit{motivator})이고, shame은 자기 변화 가능성을 좀먹는(\textit{corrodes}) 파괴적 감정이다. ③ ``\textbf{Guilt Drives Positive Change; Shame Tears Us Down}''은 이 대조를 정확히 요약한다. 따라서 정답은 ③.
\par\vspace{3pt}
\textbf{오답 소거:} ① shame이 더 강한 동기라고 역전시킴 — 본문과 반대. ② guilt와 shame이 \textit{함께} 작용한다고 암시 — 본문은 대조를 강조. ④ 사과의 역할에만 초점 — 글의 핵심 주제(대조)를 포괄하지 못함. ⑤ 둘을 동등하게 수용하라고 권고 — 본문은 guilt를 지지하고 shame을 경계함.

\sol{3. 요약문 (A)(B)}{③}
요약문: \textit{``Unlike shame, which (A) our belief in self-improvement, guilt tends to (B) positive behavioral change by pushing us to apologize and make amends.''}
\par\vspace{3pt}
(A): 지문의 ``\textit{shame \textbf{corrodes} the part of us that believes we can change}''에서 shame은 자기 개선에 대한 믿음을 \textbf{약화시킨다} → \textbf{undermines}가 적절. \textit{reinforces}(강화하다)는 반대 의미이므로 오답.
\par\vspace{2pt}
(B): 지문의 ``\textit{guilt is most often the \textbf{motivator}}''에서 guilt는 긍정적 행동 변화를 \textbf{이끈다} → \textbf{drive}가 적절. \textit{prevent}(방해하다)는 본문의 내용과 정반대이므로 오답.
\par\vspace{2pt}
따라서 (A) undermines / (B) drive → 정답 ③.

\vspace{8pt}
\noindent{\sffamily\bfseries 4. 서술형 해석}
\par\textbf{모범답안}: 사실, 내 연구에서 나는 수치심이 우리가 변화하고 더 나아질 수 있다고 믿는 우리 내면의 부분을 약화시킨다는 것을 발견했다.
\par\textbf{채점 기준}: \textit{shame corrodes}를 ``수치심이 약화시킨다/좀먹는다''로, \textit{the part of us that believes...}를 ``우리가 변화하고 더 나아질 수 있다고 믿는 부분''으로 옮기면 정답.

\end{document}
